%!TEX root = ../main.tex

\section{Theory}\label{sec:theory}

This section develops the theoretical framework.  We begin with the model
setup (Section~\ref{sec:setup}), derive the accumulation option and fair
premium (Section~\ref{sec:fair_premium}), analyse reflexivity and multiple
equilibria (Section~\ref{sec:multiple_equilibria}), reinterpret structural
indicators as option Greeks (Section~\ref{sec:greeks}), compare MSTR equity
to a BTC ETF (Section~\ref{sec:etf_comparison}), derive optimal capital
structure rules (Section~\ref{sec:optimal_capital}), and model market impact
and systemic risk (Section~\ref{sec:market_impact_theory}).

%=============================================================================
\subsection{Model Setup}\label{sec:setup}
%=============================================================================

We work on a filtered probability space
$(\Omega,\mathcal{F},(\mathcal{F}_t)_{t\ge 0},\mathbb{P})$ with the
following state variables:
\begin{itemize}
  \item $S_t$: BTC spot price;
  \item $H_t$: BTC holdings (in BTC);
  \item $D_t$: outstanding debt (face value);
  \item $N_t^{\mathrm{sh}}$: shares outstanding;
  \item $B_t = H_t / N_t^{\mathrm{sh}}$: BTC per share;
  \item $\mathrm{NAV}_t = (H_t S_t - D_t - P_t^{\mathrm{liq}})^+$: net asset
        value, where $P_t^{\mathrm{liq}}$ is the preferred stock liquidation
        value;
  \item $E_t = e^{\pi_t}\,\mathrm{NAV}_t$: equity market capitalisation;
  \item $\pi_t = \log(E_t / \mathrm{NAV}_t)$: log-premium to NAV.
\end{itemize}

We collect the observable state into the vector
\begin{equation}
  \label{eq:state_vector}
  \Phi_t = (S_t,\, H_t,\, D_t,\, N_t^{\mathrm{sh}},\, \sigma_S,\, \kappa,\, \alpha),
\end{equation}
where $\sigma_S$ is BTC volatility and $\alpha$ and $\kappa$ are behavioural
parameters defined below.

\begin{assumption}[Operating regime]
  \label{ass:operating_regime}
  We assume $A_t = H_t S_t > D_t + P_t^{\mathrm{liq}}$ so that
  $\mathrm{NAV}_t > 0$ and the log-premium $\pi_t$ is well-defined.
  The distress regime $A_t \le D_t + P_t^{\mathrm{liq}}$ is treated
  in Section~\ref{sec:death_spiral}.
\end{assumption}

\begin{assumption}[Equity issuance technology]
  \label{ass:issuance}
  The firm can issue new shares at-the-market at the prevailing equity price
  $P_t$ and immediately deploy the proceeds to purchase BTC at price $S_t$.
  Issuance is frictionless up to a market-impact cost parameterised by
  $\eta \ge 0$.
\end{assumption}

%=============================================================================
\subsection{The Accumulation Option and Fair Premium}\label{sec:fair_premium}
%=============================================================================

%-----------------------------------------------------------------------------
\subsubsection{Endogenous Premium System}\label{sec:coupled_system}
%-----------------------------------------------------------------------------

We replace an exogenous premium process with a \emph{coupled} system.

\begin{definition}[Endogenous premium system]
  \label{def:coupled_sdes}
  The joint dynamics of $(S_t, \pi_t, H_t)$ under $\mathbb{P}$ are:
  \begin{align}
    \frac{dS_t}{S_{t^-}} &= \mu_S\,dt + \sigma_S\,dW_t^S
      + \eta\,\frac{dH_t}{H_{t^-}},
    \label{eq:btc_impact} \\[6pt]
    d\pi_t &= \kappa\bigl(\pi^\star(\Phi_t) - \pi_t\bigr)\,dt
      + \sigma_\pi\,dW_t^\pi,
    \label{eq:premium_endo} \\[6pt]
    dH_t &= \alpha\,\pi_t^+\,dt + dH_t^{\mathrm{jump}},
    \label{eq:holdings_endo}
  \end{align}
  where $\mu_S$ is the BTC drift, $\eta \ge 0$ is market-impact elasticity,
  $\pi^\star(\Phi_t)$ is the endogenous fair premium derived below,
  $\alpha > 0$ governs continuous accumulation intensity,
  $\pi_t^+ = \max(\pi_t,0)$,
  $dH_t^{\mathrm{jump}} = \sum_k \Delta H_{\tau_k}\,\delta(t - \tau_k)$
  captures discrete financing events with Poisson arrival rate $\lambda_M$,
  and $\mathrm{corr}(dW_t^S, dW_t^\pi) = \rho$.
\end{definition}

\begin{remark}
  Setting $\pi^\star(\Phi_t) \equiv \theta$ (constant) recovers the exogenous
  OU premium process of the baseline model.
\end{remark}

The key novelty is that $\pi^\star$ depends on $\Phi_t$, creating a
\emph{feedback loop}:
\begin{equation}
  \label{eq:feedback_loop}
  \pi_t \uparrow \;\Longrightarrow\;
  H_t \uparrow \;\Longrightarrow\;
  B_t \uparrow \;\Longrightarrow\;
  \pi^\star \uparrow \;\Longrightarrow\;
  \pi_t \uparrow \;\cdots
\end{equation}
This is the \emph{reflexive flywheel} that distinguishes a BTC treasury
vehicle from a passive fund.

%-----------------------------------------------------------------------------
\subsubsection{Accretive Issuance: Single-Round Mechanics}
\label{sec:single_round}
%-----------------------------------------------------------------------------

\begin{definition}[Accretion from a single issuance]
  \label{def:single_accretion}
  Suppose the firm has $H$ BTC, $N$ shares, and trades at premium $\pi > 0$.
  It issues $\delta N$ new shares at market price
  $P = e^\pi \cdot (HS - D)/N$, raising proceeds to purchase
  $\delta H = e^\pi(HS-D)\delta N / (NS)$ BTC.  The change in BTC-per-share
  is
  \begin{equation}
    \label{eq:accretion_single}
    \frac{\Delta B}{B}
    = \frac{(e^\pi - 1)\,(1 - D/(HS))}{1 + \delta N / N}\cdot
      \frac{\delta N}{N}.
  \end{equation}
\end{definition}

For small issuance ($\delta N / N \ll 1$):
\begin{equation}
  \label{eq:accretion_approx}
  \frac{\Delta B}{B} \approx (e^\pi - 1)\,\frac{\delta N}{N}.
\end{equation}
Each round at $\pi > 0$ is BTC-per-share accretive.  At $\pi < 0$,
issuance is dilutive, so the firm optimally refrains: this is the
option-like payoff.

%-----------------------------------------------------------------------------
\subsubsection{Perpetual Accumulation Option Value}
\label{sec:vacc}
%-----------------------------------------------------------------------------

\begin{definition}[Accumulation option value]
  \label{def:vacc}
  The \emph{accumulation option value} is the present value of all future
  BTC-per-share accretion from optimally timed issuance:
  \begin{equation}
    \label{eq:vacc_def}
    V_{\mathrm{acc}}(\Phi_t)
    = \mathrm{NAV}_t \cdot
      \mathbb{E}^{\mathbb{P}}\!\left[
        \int_t^\infty e^{-r(u-t)}\,
        g(\pi_u,\Phi_u)\,du
        \;\Big|\;\mathcal{F}_t
      \right],
  \end{equation}
  where $r > 0$ is the discount rate and the accretion rate function is
  \begin{equation}
    \label{eq:accretion_rate}
    g(\pi,\Phi)
    = \alpha_N\,(e^\pi - 1)^+\,\biggl(1 - \frac{D}{HS}\biggr)^{\!+},
  \end{equation}
  with $\alpha_N > 0$ the issuance intensity (shares issued per unit time
  as a fraction of $N$).
\end{definition}

The $(e^\pi - 1)^+$ term captures the option-like feature: the firm issues
only when $\pi > 0$.

\begin{proposition}[Fair premium]
  \label{prop:fair_premium}
  Under the endogenous premium framework, the fair log-premium satisfies
  \begin{equation}
    \label{eq:pi_star_def}
    \boxed{
      \pi^\star(\Phi_t)
      = \log\!\left(1 + \frac{V_{\mathrm{acc}}(\Phi_t)}{\mathrm{NAV}_t}\right).
    }
  \end{equation}
\end{proposition}

\begin{proof}
  In equilibrium, $E_t = \mathrm{NAV}_t + V_{\mathrm{acc}}(\Phi_t)$.
  Since $E_t = e^{\pi_t}\,\mathrm{NAV}_t$, setting $\pi_t = \pi^\star$:
  \[
    e^{\pi^\star}\,\mathrm{NAV}_t = \mathrm{NAV}_t + V_{\mathrm{acc}},
  \]
  which gives~\eqref{eq:pi_star_def}.
\end{proof}

\begin{remark}
  $\pi^\star > 0$ if and only if $V_{\mathrm{acc}} > 0$.  A passive BTC
  holding company with no issuance capability ($\alpha_N = 0$) has
  $V_{\mathrm{acc}} = 0$ and thus $\pi^\star = 0$: it should trade at NAV.
\end{remark}

%-----------------------------------------------------------------------------
\subsubsection{Steady-State Closed-Form Approximation}
\label{sec:steady_state}
%-----------------------------------------------------------------------------

\begin{assumption}[Steady-state approximation]
  \label{ass:steady_state}
  The premium $\pi_t$ fluctuates around $\pi^\star$ as a stationary OU
  process with speed $\kappa$ and volatility $\sigma_\pi$.  The state
  variables $(S_t, H_t, D_t)$ evolve slowly relative to premium
  mean-reversion.
\end{assumption}

Under this assumption, $\pi \sim \mathcal{N}(\pi^\star, (\sigma_\pi^{\mathrm{ss}})^2)$
where $\sigma_\pi^{\mathrm{ss}} = \sigma_\pi / \sqrt{2\kappa}$.

\begin{proposition}[Steady-state fair premium]
  \label{prop:steady_state}
  With leverage ratio $\ell = D/(HS)$ and $L = (1-\ell)^+$, the fair
  premium satisfies the fixed-point equation
  \begin{equation}
    \label{eq:pi_star_fp}
    \pi^\star
    = \log\!\left(
        1 + \frac{\alpha_N L}{r}\,
        \mathbb{E}\bigl[(e^{\pi} - 1)^+\bigr]
      \right),
  \end{equation}
  where $\pi \sim \mathcal{N}(\pi^\star,\, (\sigma_\pi^{\mathrm{ss}})^2)$.
  Evaluating the expectation:
  \begin{equation}
    \label{eq:pi_star_closed}
    \boxed{
      \pi^\star
      = \log\!\left(
          1 + \frac{\alpha_N L}{r}\,\Bigl[
            e^{\pi^\star + (\sigma_\pi^{\mathrm{ss}})^2/2}\,
            \Phi\!\left(\frac{\pi^\star}{\sigma_\pi^{\mathrm{ss}}}
              + \sigma_\pi^{\mathrm{ss}}\right)
            - \Phi\!\left(\frac{\pi^\star}{\sigma_\pi^{\mathrm{ss}}}\right)
          \Bigr]
        \right),
    }
  \end{equation}
  where $\Phi(\cdot)$ is the standard normal CDF.
\end{proposition}

\begin{proof}
  At steady state, $V_{\mathrm{acc}}/\mathrm{NAV} = (\alpha_N L / r)\,
  \mathbb{E}[(e^\pi - 1)^+]$.  The expectation with
  $\pi \sim \mathcal{N}(m, v^2)$ is the Black--Scholes formula at strike
  $K=1$:
  \[
    \mathbb{E}[(e^\pi - 1)^+]
    = e^{m + v^2/2}\,\Phi\!\left(\frac{m}{v} + v\right)
      - \Phi\!\left(\frac{m}{v}\right),
  \]
  with $m = \pi^\star$, $v = \sigma_\pi^{\mathrm{ss}}$.
  Substituting into~\eqref{eq:pi_star_def} yields~\eqref{eq:pi_star_closed}.
\end{proof}

%-----------------------------------------------------------------------------
\subsubsection{Comparative Statics}\label{sec:comparative_statics}
%-----------------------------------------------------------------------------

\begin{proposition}[Comparative statics of the fair premium]
  \label{prop:comp_statics}
  Let $F(\pi^\star;\,\sigma_S,\kappa,\alpha_N,\ell,r)$ denote the
  right-hand side of~\eqref{eq:pi_star_closed}.  At any interior fixed point
  $\pi^\star > 0$ with $\partial F/\partial \pi^\star < 1$ (stable
  equilibrium):
  \begin{enumerate}
    \item $\partial \pi^\star / \partial \sigma_S > 0$:
      higher BTC volatility increases the option value of future accretive
      issuance.
    \item $\partial \pi^\star / \partial \kappa < 0$:
      faster mean-reversion reduces stationary variance and hence option
      value.
    \item $\partial \pi^\star / \partial \alpha_N > 0$:
      greater issuance ability increases option value (the IBGR channel).
    \item $\partial \pi^\star / \partial \ell < 0$ for $\ell < 1$:
      higher leverage reduces accretion efficiency.
    \item $\partial \pi^\star / \partial r < 0$:
      higher discount rates reduce present value of future accretion.
  \end{enumerate}
\end{proposition}

\begin{proof}
  By the implicit function theorem at a fixed point with $F_{\pi^\star} < 1$:
  $\partial \pi^\star / \partial p_i = (\partial F/\partial p_i) /
  (1 - \partial F / \partial \pi^\star)$.
  The denominator is positive at a stable equilibrium.  The signs of the
  numerators follow from direct differentiation
  of~\eqref{eq:pi_star_closed}: (1)~the Black--Scholes call value has
  positive vega; (2)~$\sigma_\pi^{\mathrm{ss}}$ is decreasing in $\kappa$;
  (3)~$F$ is increasing in $\alpha_N$; (4)~$L = (1-\ell)^+$ is decreasing
  in $\ell$; (5)~$r$ appears in the denominator.
\end{proof}

\begin{corollary}[Conditions for positive fair premium]
  \label{cor:sign_pistar}
  $\pi^\star > 0$ if and only if $\alpha_N > 0$, $L > 0$ (i.e.\
  $\ell < 1$), and $\sigma_\pi^{\mathrm{ss}} > 0$.
  A passive holding company ($\alpha_N = 0$) or one at maximum leverage
  ($\ell \ge 1$) has $\pi^\star = 0$.
\end{corollary}

%=============================================================================
\subsection{Reflexivity and Multiple Equilibria}
\label{sec:multiple_equilibria}
%=============================================================================

%-----------------------------------------------------------------------------
\subsubsection{Feedback Gain}\label{sec:feedback_gain}
%-----------------------------------------------------------------------------

\begin{definition}[Feedback gain]
  \label{def:feedback_gain}
  The feedback gain at equilibrium $\pi^\star$ is
  \begin{equation}
    \label{eq:feedback_gain}
    G(\pi^\star)
    = \frac{\partial F}{\partial \pi^\star},
  \end{equation}
  measuring the marginal increase in the fair premium from a unit increase
  in the current premium level.
\end{definition}

\begin{proposition}[Effective mean-reversion]
  \label{prop:effective_kappa}
  Near a stable equilibrium with feedback gain $G < 1$, the effective
  mean-reversion speed is
  \begin{equation}
    \label{eq:kappa_eff}
    \boxed{
      \kappa_{\mathrm{eff}} = \kappa\,(1 - G).
    }
  \end{equation}
\end{proposition}

\begin{proof}
  Linearise $\pi^\star(\Phi_t)$ around steady state.  When $\pi_t$ increases
  by $\delta\pi$, the target shifts by $G\,\delta\pi$.  The restoring drift
  becomes $-\kappa(1 - G)\,\delta\pi$.
\end{proof}

Three regimes emerge: (i)~$G < 1$: stable, premium mean-reverts more slowly
than the exogenous model; (ii)~$G = 1$: critical, mean-reversion vanishes;
(iii)~$G > 1$: unstable, perturbations self-amplify.

%-----------------------------------------------------------------------------
\subsubsection{Fixed-Point Analysis and Equilibrium Multiplicity}
\label{sec:fixed_point}
%-----------------------------------------------------------------------------

\begin{theorem}[Existence and multiplicity of equilibria]
  \label{thm:three_equilibria}
  Consider the fixed-point equation $\pi^\star = F(\pi^\star)$
  from~\eqref{eq:pi_star_closed} with parameters
  $(\alpha_N, L, r, \sigma_\pi^{\mathrm{ss}})$ all strictly positive.
  Define the feedback strength
  \begin{equation}
    \label{eq:feedback_strength}
    \Gamma = \frac{\alpha_N L}{r}.
  \end{equation}
  Then:
  \begin{enumerate}
    \item \textbf{Weak feedback:} If $\Gamma$ is sufficiently small, there
      exists a unique fixed point $\pi^\star_{\mathrm{low}} > 0$ (stable,
      $G < 1$).
    \item \textbf{Strong feedback:} If $\Gamma$ is sufficiently large,
      three fixed points exist:
      $\pi^\star_{\mathrm{low}}$ (stable),
      $\pi^\star_{\mathrm{mid}}$ (unstable tipping point, $G > 1$), and
      $\pi^\star_{\mathrm{high}}$ (stable flywheel).
    \item \textbf{Bifurcation:} The transition occurs at a critical
      $\Gamma^*$ via saddle-node bifurcation where
      $\pi^\star_{\mathrm{mid}}$ and $\pi^\star_{\mathrm{high}}$ collide.
  \end{enumerate}
\end{theorem}

\begin{proof}
  Define $h(\pi^\star) = F(\pi^\star) - \pi^\star$.  We have $h(0) > 0$
  (since $\sigma_\pi^{\mathrm{ss}} > 0$ implies positive option value at
  $\pi^\star = 0$).  For small $\Gamma$, $F' < 1$ everywhere, so $h$ is
  strictly decreasing with a unique zero.  For large $\Gamma$, the
  exponential growth of the option value creates a region where $F' > 1$,
  producing the S-shaped curve with three crossings.  At the critical
  $\Gamma^*$, $h(\pi^\star_c) = 0$ and $h'(\pi^\star_c) = 0$
  simultaneously---the standard saddle-node bifurcation.  Stability follows
  from $G < 1$ at the low and high equilibria and $G > 1$ at the middle.
\end{proof}

\begin{remark}[Economic interpretation]
  \label{rem:three_equil_econ}
  The \textbf{flywheel} ($\pi^\star_{\mathrm{high}}$): high premium $\to$
  aggressive issuance $\to$ rapid accretion $\to$ justified high premium.
  The \textbf{tipping point} ($\pi^\star_{\mathrm{mid}}$): unstable
  separator between basins of attraction.
  The \textbf{spiral} ($\pi^\star_{\mathrm{low}}$): low premium $\to$
  minimal issuance $\to$ negligible accretion $\to$ justified low premium.
\end{remark}

%-----------------------------------------------------------------------------
\subsubsection{Hysteresis}\label{sec:hysteresis}
%-----------------------------------------------------------------------------

\begin{proposition}[Hysteresis in premium transitions]
  \label{prop:hysteresis}
  Consider a parameter path where $\Gamma$ increases through $\Gamma^*$
  and then decreases.
  \begin{enumerate}
    \item As $\Gamma$ increases past $\Gamma^*$, the system at
      $\pi^\star_{\mathrm{low}}$ does not jump to
      $\pi^\star_{\mathrm{high}}$; it remains at the low equilibrium until
      a large shock exceeds the tipping point.
    \item If at $\pi^\star_{\mathrm{high}}$ and $\Gamma$ decreases below
      $\Gamma^*$, the flywheel equilibrium vanishes and the premium
      collapses discontinuously to $\pi^\star_{\mathrm{low}}$.
    \item The path exhibits hysteresis: the equilibrium depends on the
      direction of approach.
  \end{enumerate}
\end{proposition}

\begin{proof}
  Standard saddle-node bifurcation: when the stable (flywheel) and unstable
  (tipping) branches annihilate, no nearby attractor exists, forcing a
  discontinuous transition.
\end{proof}

%-----------------------------------------------------------------------------
\subsubsection{Death Spiral Mechanism}\label{sec:death_spiral}
%-----------------------------------------------------------------------------

\begin{theorem}[Death spiral]
  \label{thm:death_spiral}
  If $\pi_t < 0$ (discount to NAV):
  \begin{enumerate}
    \item $V_{\mathrm{acc}} \approx 0$ since issuance at a discount is
      dilutive and the firm refrains.
    \item The fair premium collapses to $\pi^\star \approx 0$.
    \item The feedback loop reverses:
      $\pi_t < 0 \Rightarrow \alpha_N^{\mathrm{eff}} = 0 \Rightarrow
      V_{\mathrm{acc}} \downarrow \Rightarrow \pi^\star \downarrow
      \Rightarrow \pi_t \downarrow$.
    \item A death spiral occurs when
      \begin{equation}
        \label{eq:spiral_condition}
        e^{\pi_t} < \frac{D_t}{H_t S_t}
        \quad\Longleftrightarrow\quad
        E_t < D_t,
      \end{equation}
      i.e.\ the market values equity below the debt burden.
  \end{enumerate}
\end{theorem}

\begin{proof}
  When $\pi_t \le 0$, $(e^\pi - 1)^+ = 0$, so $V_{\mathrm{acc}} = 0$ and
  $\pi^\star = 0$.  The natural OU restoring drift is
  $\kappa(0 - \pi_t) > 0$, tending toward recovery.  However, if forced
  dilutive issuance (to meet debt service) exceeds natural mean-reversion,
  the vicious cycle takes hold when $E_t < D_t$.
\end{proof}

%-----------------------------------------------------------------------------
\subsubsection{Mispricing Metric and Z-Score Bands}
\label{sec:mispricing}
%-----------------------------------------------------------------------------

\begin{definition}[Premium mispricing]
  \label{def:mispricing}
  The mispricing is $\Delta_t = \pi_t - \pi^\star(\Phi_t)$.
  Positive $\Delta_t$: overvaluation; negative: undervaluation.
\end{definition}

\begin{proposition}[Dynamics of the mispricing]
  \label{prop:delta_dynamics}
  Under the slow-variation assumption, $\Delta_t$ satisfies
  \begin{equation}
    \label{eq:delta_ou}
    d\Delta_t \approx -\kappa\,\Delta_t\,dt
      + \sigma_\pi\,dW_t^\pi
      - d\pi^\star(\Phi_t),
  \end{equation}
  with mispricing mean-reversion speed
  \begin{equation}
    \label{eq:kappa_delta}
    \boxed{
      \kappa_\Delta = \kappa > \kappa_{\mathrm{eff}}.
    }
  \end{equation}
  The mispricing reverts faster than the premium itself.
\end{proposition}

\begin{proof}
  Write $\pi_t = \pi^\star + \Delta_t$.  From~\eqref{eq:premium_endo},
  $d\pi_t = -\kappa\Delta_t\,dt + \sigma_\pi dW_t^\pi$.  Subtracting
  $d\pi^\star$: $d\Delta_t = -\kappa\Delta_t\,dt + \sigma_\pi dW_t^\pi
  - d\pi^\star$.  Since $\kappa > \kappa(1-G) = \kappa_{\mathrm{eff}}$
  for $G > 0$, the mispricing reverts faster.
\end{proof}

\begin{definition}[Mispricing Z-score]
  \label{def:zscore}
  \begin{equation}
    \label{eq:zscore}
    Z_t^{\Delta} = \frac{\Delta_t}{\sigma_\pi / \sqrt{2\kappa}}.
  \end{equation}
\end{definition}

Under the OU approximation, $Z_t^\Delta$ is approximately standard normal,
yielding classification bands: $|Z| < 1$ (fair value), $1$--$2$ (moderate
mispricing), $> 2$ (extreme mispricing, mean-reversion expected).

\begin{proposition}[Half-life of mispricing]
  \label{prop:halflife}
  $t_{1/2}^{\Delta} = \ln 2 / \kappa$, shorter than the premium half-life
  $\ln 2 / \kappa_{\mathrm{eff}}$ by the factor $(1-G)$.
\end{proposition}

%=============================================================================
\subsection{Structural Indicators as Option Greeks}\label{sec:greeks}
%=============================================================================

The structural relation $E_t = e^{\pi_t}(A_t - D_t - P_t^{\mathrm{liq}})^+$
endows MSTR equity with an option-like payoff on Bitcoin.  If $\pi_t$ encodes
the accumulation option value, then the structural indicators are precisely
the option Greeks of this embedded claim.

\begin{definition}[Accumulation option value]
  \label{def:acc_option_greeks}
  $V_{\mathrm{acc}}(t) = E_t - \mathrm{NAV}_t = (e^{\pi_t} - 1)(A_t - D_t - P_t^{\mathrm{liq}})^+$.
  When $\pi_t > 0$ the option is in the money; when $\pi_t \le 0$ it is
  worthless.
\end{definition}

%-----------------------------------------------------------------------------
\subsubsection{ILE as Delta}
%-----------------------------------------------------------------------------

\begin{proposition}[ILE $\equiv$ Leverage Delta]\label{prop:ile_delta}
  With $\pi_t$ held fixed:
  \begin{equation}
    \Delta_{\mathrm{lev}}(t)
    = \frac{\partial \log E_t}{\partial \log S_t}\bigg|_{\pi_t}
    = \frac{A_t}{A_t - D_t - P_t^{\mathrm{liq}}}
    = \mathrm{ILE}_t.
  \end{equation}
\end{proposition}

\begin{proof}
  $\log E_t = \pi_t + \log(H_tS_t - D_t - P_t^{\mathrm{liq}})$.
  Differentiating w.r.t.\ $\log S_t$ with $H_t, D_t, P_t^{\mathrm{liq}}$
  fixed gives $H_tS_t / (H_tS_t - D_t - P_t^{\mathrm{liq}}) = \mathrm{ILE}_t$.
\end{proof}

With calibrated values, $\mathrm{ILE}_0 = 1.49$: a 1\% BTC move produces a
1.49\% equity move from leverage alone.

%-----------------------------------------------------------------------------
\subsubsection{TEE as Adjusted Delta}
%-----------------------------------------------------------------------------

\begin{proposition}[TEE $\equiv$ Total Delta]\label{prop:tee_delta}
  Let $\gamma_{\pi S}$ denote the sensitivity of the log-premium to BTC
  log-returns.  Then
  \begin{equation}
    \Delta_{\mathrm{tot}}(t)
    = \frac{d\log E_t}{d\log S_t}
    = \gamma_{\pi S} + \mathrm{ILE}_t
    = \mathrm{TEE}_t.
  \end{equation}
\end{proposition}

\begin{proof}
  Total derivative: $d\log E_t / d\log S_t = d\pi_t/d\log S_t + A_t/(A_t - D_t - P_t^{\mathrm{liq}})$.
\end{proof}

At current calibration, $\mathrm{TEE}_0 = -3.60 + 1.49 = -2.11$: net equity
moves inversely to BTC, reflecting a regime where premium compression
dominates leverage.

%-----------------------------------------------------------------------------
\subsubsection{PMRI as Theta}
%-----------------------------------------------------------------------------

\begin{proposition}[PMRI $\equiv$ Premium Theta]\label{prop:pmri_theta}
  Under OU dynamics $d\pi_t = \kappa(\theta - \pi_t)\,dt + \sigma_\pi dW_t^\pi$,
  the standardised expected premium decay is
  \begin{equation}
    \Theta_\pi(t)
    = \frac{\pi_t - \theta}{\sigma_\pi / \sqrt{2\kappa}}
    = \mathrm{PMRI}_t.
  \end{equation}
\end{proposition}

With $\kappa = 5.21$, $\theta = 1.08$, $\sigma_\pi = 3.79$, and
$\pi_0 = 0.24$:
\[
  \sigma_{\mathrm{stat}} = \frac{3.79}{\sqrt{2 \times 5.21}} = 1.175,
  \quad
  \mathrm{PMRI}_0 = \frac{0.24 - 1.08}{1.175} = -0.71.
\]
The negative value implies the premium is below its long-run mean,
conferring positive expected carry.

%-----------------------------------------------------------------------------
\subsubsection{IBGR as Moneyness}
%-----------------------------------------------------------------------------

\begin{proposition}[IBGR $\equiv$ Moneyness]\label{prop:ibgr_moneyness}
  The implied BTC growth rate
  \begin{equation}
    \mathrm{IBGR}_t = \mu_H(t)
    = \frac{1}{H_t}\bigl(\alpha\,\pi_t^+ + \lambda_M\,\mathbb{E}[\Delta H]\bigr)
  \end{equation}
  measures how deeply in the money the accumulation option is.
\end{proposition}

With current calibration, $\mathrm{IBGR}_0 = (0 + 19.05 \times 7{,}044) /
761{,}068 = 17.6\%$ annualised.

%-----------------------------------------------------------------------------
\subsubsection{IFRD as Exercise Cost}
%-----------------------------------------------------------------------------

\begin{proposition}[IFRD $\equiv$ Exercise Cost]\label{prop:ifrd_strike}
  The implied funding requirement $G_T = \int_0^T S_u\,dH_u$ is the
  cumulative exercise cost of the accumulation option.  From simulation:
  $\mathbb{E}[G_{3\mathrm{y}}] = \$35.3$B, $G_{3\mathrm{y}}^{95\%} = \$71.8$B.
\end{proposition}

%-----------------------------------------------------------------------------
\subsubsection{Summary and New Greeks}
%-----------------------------------------------------------------------------

\begin{table}[H]
\centering
\caption{Structural indicators as option Greeks.}
\label{tab:greek-map}
\begin{tabular}{@{}llll@{}}
\toprule
\textbf{Indicator} & \textbf{Greek} & \textbf{Formula} & \textbf{Calibrated} \\
\midrule
ILE $= A/(A-D-P^{\mathrm{liq}})$ & Delta ($\Delta_{\mathrm{lev}}$)
  & $\partial\log E / \partial\log S\big|_\pi$
  & $1.49$ \\[4pt]
TEE $= \gamma_{\pi S} + \text{ILE}$ & Adjusted Delta ($\Delta_{\mathrm{tot}}$)
  & $d\log E / d\log S$
  & $-2.11$ \\[4pt]
PMRI $= (\pi-\theta)/\sigma_{\mathrm{stat}}$ & Theta ($\Theta_\pi$)
  & Premium carry cost
  & $-0.71$ \\[4pt]
IBGR $= \mu_H$ & Moneyness
  & $(\alpha\pi^+ + \lambda_M\mathbb{E}[\Delta H])/H$
  & $17.6\%$ \\[4pt]
IFRD $= G_T$ & Exercise cost
  & $\int S\,dH$
  & \$35.3B (mean) \\
\bottomrule
\end{tabular}
\end{table}

We introduce two additional Greeks absent from the original indicator set:

\begin{definition}[Premium Vega]
  $\mathcal{V}_\pi = \partial\pi^*/\partial\sigma_S$.
\end{definition}

\begin{proposition}[Positive Vega]\label{prop:vega}
  The accumulation option has $\mathcal{V}_\pi > 0$: higher BTC volatility
  increases the fair premium through option convexity and more frequent
  flywheel episodes.
\end{proposition}

\begin{proof}
  The equity payoff $(A_t - D_t)^+$ is convex in $S_t$ (Jensen's inequality),
  and higher $\sigma_S$ generates more accumulation opportunities through the
  drift $\alpha\pi_t^+$.
\end{proof}

\begin{definition}[Premium Rho]
  $\mathcal{R}_\pi = \partial\pi^*/\partial r$.
\end{definition}

\begin{proposition}[Negative Rho]\label{prop:rho_negative}
  $\mathcal{R}_\pi < 0$: higher interest rates decrease the fair premium
  through increased debt servicing cost, higher financing cost, and higher
  discount rate on future accretion.
\end{proposition}

%=============================================================================
\subsection{MSTR vs.\ BTC ETF: Premium Decomposition}
\label{sec:etf_comparison}
%=============================================================================

The launch of spot Bitcoin ETFs created a pure delta-one BTC exposure at
NAV.  Why does MSTR trade at a premium?

\begin{definition}[BTC ETF]\label{def:etf}
  A spot BTC ETF satisfies $E_t^{\mathrm{ETF}} = H_t^{\mathrm{ETF}}S_t$,
  with $\pi^{\mathrm{ETF}} = 0$, no debt, no leverage ($\mathrm{ILE} = 1$),
  no accumulation option, and $\Delta_{\mathrm{tot}}^{\mathrm{ETF}} = 1$.
\end{definition}

\begin{theorem}[Three-component premium decomposition]
  \label{thm:decomposition}
  The MSTR log-premium decomposes as
  \begin{equation}
    \pi_t
    = \underbrace{\pi_{\mathrm{lev}}(t)}_{\text{leverage}}
    + \underbrace{\pi_{\mathrm{acc}}(t)}_{\text{accumulation}}
    + \underbrace{\pi_{\mathrm{noise}}(t)}_{\text{residual}},
  \end{equation}
  where $\pi_{\mathrm{lev}} = \log(\mathrm{ILE}_t)$ and
  $\pi_{\mathrm{acc}} = \log(1 + V_{\mathrm{acc}}/\mathrm{NAV}_t)$.
\end{theorem}

\begin{proof}
  From $E_t = \mathrm{NAV}_t + V_{\mathrm{acc}} + \varepsilon_t$,
  taking logs: $\pi_t = \log(E_t/\mathrm{NAV}_t)$.  The leverage premium
  $\pi_{\mathrm{lev}} = \log(A_t/(A_t - D_t - P_t^{\mathrm{liq}}))$
  captures the excess sensitivity from the levered structure, the
  accumulation premium captures option value, and the residual absorbs
  sentiment and noise.
\end{proof}

At current calibration: $\pi_{\mathrm{lev}} = \log(1.488) = 0.398$ and
$\pi_0 - \pi_{\mathrm{lev}} = 0.242 - 0.398 = -0.156$, implying the
current premium does not fully compensate for leverage risk---the
accumulation option is priced at a discount.

\begin{proposition}[MSTR dominance conditions]\label{prop:mstr_dominates}
  MSTR offers superior risk-adjusted BTC exposure when:
  (i)~$\pi_{\mathrm{acc}} > 0$ (accumulation option in the money),
  (ii)~$\mathrm{PMRI}_t < 0$ (premium underpriced), and
  (iii)~$\mathrm{ILE}_t < \bar\beta$ (sustainable leverage).
\end{proposition}

\begin{proof}
  The expected MSTR excess return over the ETF is
  $\mathbb{E}[\Delta\pi_t] + (\mathrm{ILE}_t - 1)\mathbb{E}[\Delta\log S_t]$.
  Under conditions (i)--(iii), the first term is positive (premium expansion),
  the second provides leverage upside, and leverage risk is bounded.
\end{proof}

\begin{proposition}[ETF dominance conditions]\label{prop:etf_dominates}
  The ETF dominates when any of: $\mathrm{PMRI}_t > 0$ (overpriced premium),
  $\mathrm{ILE}_t > \bar\beta$ (excessive leverage),
  $\mathrm{IBGR}_t \le 0$ (stalled flywheel), or deteriorating dividend
  coverage.
\end{proposition}

\begin{corollary}[Premium regime map]\label{cor:regime_map}
  \begin{center}
  \begin{tabular}{@{}lcc@{}}
  \toprule
  & $\mathrm{IBGR} > 0$ & $\mathrm{IBGR} \le 0$ \\
  \midrule
  $\mathrm{PMRI} < 0$ (underpriced)
    & \textbf{MSTR dominates} & Transition / distress \\
  $\mathrm{PMRI} > 0$ (overpriced)
    & Hold / hedge & \textbf{ETF dominates} \\
  \bottomrule
  \end{tabular}
  \end{center}
  At current calibration ($\mathrm{IBGR} = 17.6\%$,
  $\mathrm{PMRI} = -0.71$), MSTR dominates.
\end{corollary}

\begin{table}[H]
\centering
\caption{MSTR vs.\ spot BTC ETF: structural comparison.}
\label{tab:mstr-vs-etf}
\begin{tabular}{@{}lcc@{}}
\toprule
\textbf{Metric} & \textbf{MSTR} & \textbf{BTC ETF} \\
\midrule
BTC Delta (ILE) & 1.49 & 1.00 \\
Total Delta (TEE) & $-2.11$ & 1.00 \\
Premium ($\pi$) & 0.24 (27\% over NAV) & $\approx 0$ \\
Accumulation option & In the money & None \\
IBGR (annualised) & 17.6\% & 0\% \\
Premium carry (PMRI) & $-0.71$ (positive carry) & N/A \\
Leverage risk & Moderate (ILE $< 2$) & None \\
3-year survival & 91.9\% & 100\% \\
Preferred dividend drag & \$977M/year & None \\
Funding requirement (3y) & \$35.3B mean & None \\
\bottomrule
\end{tabular}
\end{table}

%=============================================================================
\subsection{Optimal Capital Structure and Financing}\label{sec:optimal_capital}
%=============================================================================

Strategy employs five financing channels: ATM common equity, convertible
senior notes, and three classes of perpetual preferred stock.

%-----------------------------------------------------------------------------
\subsubsection{ATM Equity}
%-----------------------------------------------------------------------------

\begin{proposition}[ATM accretion threshold]\label{prop:atm_threshold}
  ATM equity issuance is BTC-per-share accretive if and only if
  \begin{equation}
    \label{eq:atm_accretive}
    \pi_t > \log(\mathrm{ILE}_t)
    = \log\!\left(\frac{H_t S_t}{H_t S_t - D_t - P_t^{\mathrm{liq}}}\right).
  \end{equation}
  When accretive, the marginal cost of BTC via ATM is negative (value-creating).
\end{proposition}

\begin{proof}
  Each new share at price $P_t$ purchases $P_t/S_t$ BTC.  Accretive iff
  $P_t/S_t > B_t$, i.e.\ $e^{\pi_t}(A_t - D_t - P_t^{\mathrm{liq}})/(N_tS_t) > H_t/N_t$,
  giving $e^{\pi_t} > A_t/(A_t - D_t - P_t^{\mathrm{liq}}) = \mathrm{ILE}_t$.
\end{proof}

With $\mathrm{ILE}_0 = 1.488$, the threshold is $\pi_t > 0.397$.  The
current premium $\pi_0 = 0.242$ is below this threshold, indicating ATM
issuance would be mildly dilutive---consistent with Strategy's observed
shift toward preferred stock issuance.

%-----------------------------------------------------------------------------
\subsubsection{Convertible Senior Notes}
%-----------------------------------------------------------------------------

A convertible note with coupon $c_j$, conversion price $K_j$, and maturity
$T_j$ has effective cost:
\begin{equation}
  \label{eq:conv_cost}
  k_{\mathrm{conv},j} = c_j + \frac{1}{T_j - t}\,
  \mathbb{E}^{\mathbb{P}}\!\left[
    \left(\frac{D_j}{K_j}\cdot\frac{B_{T_j}}{1 + D_j/(K_jN_{T_j}^{\mathrm{sh}})}
    - \frac{D_j}{S_{T_j}}\right)^{\!+}
    \cdot \frac{S_{T_j}}{D_j}
  \right].
\end{equation}
Strategy's convertibles carry a weighted average coupon of 0.42\%, making
them the cheapest debt channel.

%-----------------------------------------------------------------------------
\subsubsection{Preferred Stock}
%-----------------------------------------------------------------------------

For non-convertible preferred (STRF at 10\%, STRD at 10\%, STRE at 10\%,
STRC at 9.6\% initial / 11.5\% current), the cost is
$k_{\mathrm{pref},j} = d_j L_j / \text{Proceeds}_j$, yielding effective
costs of 10.7\%--12.5\%.  The convertible preferred STRK (8\%, conversion at
\$1,000) has effective cost approximately 10\% with conversion currently deep
out-of-the-money.  Total annual preferred dividend: \$976.6M.

%-----------------------------------------------------------------------------
\subsubsection{Optimal Financing Rule}
%-----------------------------------------------------------------------------

\begin{proposition}[Optimal channel selection]\label{prop:optimal_channel}
  The cost-minimising financing rule is:
  \begin{equation}
    \label{eq:optimal_channel}
    c^*(\pi_t) =
    \begin{cases}
      \text{ATM equity} & \text{if } \pi_t > \pi^*_{\mathrm{ATM}},
        \quad k_{\mathrm{ATM}} < 0, \\[4pt]
      \text{Convertible bonds} & \text{if } \pi^*_{\mathrm{conv}} < \pi_t
        \le \pi^*_{\mathrm{ATM}},
        \quad k_{\mathrm{conv}} \approx 0.4\text{--}2.3\%, \\[4pt]
      \text{Preferred stock} & \text{if } \pi_t \le \pi^*_{\mathrm{conv}},
        \quad k_{\mathrm{pref}} \approx 10\text{--}12.5\%.
    \end{cases}
  \end{equation}
\end{proposition}

Strategy's observed behaviour closely matches: massive ATM and convertible
issuance in Q4~2024 ($\pi_t \approx 1.5$--$2.0$), shift to STRK/STRF in
early 2025, then heavy STRC/STRD/STRE issuance in H2~2025--Q1~2026 as the
premium declined to 0.2--0.5.

%-----------------------------------------------------------------------------
\subsubsection{WACBA and Leverage Dynamics}
%-----------------------------------------------------------------------------

\begin{definition}[WACBA]\label{def:wacba}
  The Weighted Average Cost of BTC Acquisition:
  \begin{equation}
    \mathrm{WACBA}_t = \sum_{c \in \mathcal{C}} w_c(t)\, k_c(\pi_t, S_t),
  \end{equation}
  where $w_c(t)$ is the share of BTC acquired through channel $c$.
\end{definition}

\begin{proposition}[Procyclical leverage]\label{prop:procyclical}
  Leverage dynamics are procyclical:
  \begin{enumerate}
    \item \textbf{Bull market:} ATM issuance is accretive, ILE falls toward~1.
    \item \textbf{Bear market:} Financing shifts to preferred (perpetual
      obligations), ILE rises as $A_t \to D_t + P_t^{\mathrm{liq}}$.
  \end{enumerate}
\end{proposition}

\begin{proof}
  $\partial\mathrm{ILE}_t/\partial S_t = -(D_t + P_t^{\mathrm{liq}})H_t /
  (H_tS_t - D_t - P_t^{\mathrm{liq}})^2 < 0$.
\end{proof}

\begin{corollary}[Leverage spiral condition]\label{cor:leverage_spiral}
  Instability arises when the preferred dividend burden grows faster than
  expected BTC asset income:
  \begin{equation}
    \frac{d}{dt}\!\left(\sum_j d_j L_j\right) >
    \frac{d}{dt}\!\left(\mu_S \cdot H_t S_t\right).
  \end{equation}
  At current parameters, the right-hand side (\$6.24B at $\mu_S = 11.1\%$)
  far exceeds \$976.6M in annual preferred dividends, providing substantial
  buffer.  A sustained BTC decline to \$15,000--\$20,000 would close this gap.
\end{corollary}

%=============================================================================
\subsection{BTC Price Impact and Systemic Risk}
\label{sec:market_impact_theory}
%=============================================================================

Strategy's purchases are large relative to daily volume.  We augment BTC
price dynamics with permanent price impact:
\begin{equation}
  \label{eq:price_impact}
  \frac{dS_t}{S_t} = \mu_S\,dt + \sigma_S\,dW_t^S
  + \eta\,\frac{dH_t}{V_t},
\end{equation}
where $V_t$ is daily BTC volume and $\eta > 0$ is the Kyle lambda.

Substituting holding dynamics:
\begin{equation}
  \frac{dS_t}{S_t} = \left(\mu_S + \eta\,\frac{\alpha\pi_t^+}{V_t}\right)dt
  + \sigma_S\,dW_t^S + \eta\,\frac{dH_t^{\mathrm{jump}}}{V_t}.
\end{equation}

%-----------------------------------------------------------------------------
\subsubsection{Reflexivity Amplification}
%-----------------------------------------------------------------------------

\begin{definition}[Feedback gain with market impact]
  \begin{equation}
    G_t = \eta\,\frac{\phi_t\,\beta_{\mathrm{TOT}}(t)\,E_t}{V_t\,S_t^2},
  \end{equation}
  where $\phi_t$ is the reinvestment rate and
  $\beta_{\mathrm{TOT}} = \mathrm{ILE}_t + \gamma_{\pi S}$.
\end{definition}

\begin{proposition}[Reflexivity amplification factor]\label{prop:reflexivity}
  The steady-state amplification of an exogenous BTC demand shock is
  \begin{equation}
    R_t = \frac{1}{1 - G_t},
  \end{equation}
  provided $G_t < 1$.
\end{proposition}

\begin{proof}
  A unit exogenous price increase $\delta$ generates feedback purchases of
  $G_t\delta/\eta$, producing second-round impact $G_t\delta$.  Iterating:
  $\delta(1 + G_t + G_t^2 + \cdots) = \delta/(1-G_t)$.
\end{proof}

\begin{proposition}[Instability condition]\label{prop:instability}
  The loop becomes unstable when
  $\eta\,\phi_t\,\beta_{\mathrm{TOT}}(t)\,E_t \ge V_t\,S_t^2$.
  At current calibration, $\beta_{\mathrm{TOT}} = -2.11 < 0$, so $G_t < 0$
  and $R_t < 1$: the premium compression dampens rather than amplifies price
  impact.
\end{proposition}

With market-impact augmentation, the effective feedback gain becomes:
\begin{equation}
  G_{\mathrm{aug}} = G + \eta\,\frac{\alpha_N L}{\kappa}\,\beta_{LS}(t).
\end{equation}

%-----------------------------------------------------------------------------
\subsubsection{Discrete Financing Events}
%-----------------------------------------------------------------------------

\begin{theorem}[Expected price impact of financing program]
  \label{thm:expected_impact}
  Over horizon $[0,T]$:
  \begin{equation}
    \mathbb{E}\!\left[\int_0^T \eta\,\frac{dH_u}{V_u}\right]
    = \eta\left(
      \frac{\alpha\,\mathbb{E}[\pi_t^+]}{\bar{V}}\,T
      + \lambda_M\,\frac{\mathbb{E}[\Delta H]}{\bar{V}}\,T
    \right).
  \end{equation}
  With calibrated parameters ($\lambda_M = 19.05$,
  $\mathbb{E}[\Delta H] = 7{,}044$, $\bar{V} \approx 25{,}000$ BTC/day):
  annual price contribution $\approx 5.37\eta$, or 0.54--2.68\% for
  $\eta \in [0.1, 0.5]$.
\end{theorem}

%-----------------------------------------------------------------------------
\subsubsection{Systemic Risk from Corporate Adoption}
%-----------------------------------------------------------------------------

Strategy holds $\sim$3.84\% of circulating BTC supply.

\begin{proposition}[Fire-sale externality]\label{prop:fire_sale}
  If $F$ firms each hold BTC with leverage, the aggregate fire-sale impact is
  \begin{equation}
    \frac{\Delta S^{\mathrm{fire}}}{S_t}
    = -\eta_{\mathrm{sell}}\,
    \frac{\sum_{f: S_t < S^{(f)}_{\mathrm{distress}}} \Delta H^{(f)}_{\mathrm{sell}}}{V_t}.
  \end{equation}
\end{proposition}

\begin{theorem}[Contagion cascade condition]\label{thm:contagion}
  A cascade propagates from firm $1$ to firm $j$ if
  \begin{equation}
    S_t - \eta_{\mathrm{sell}}\,S_t\,
    \frac{\sum_{i=1}^{j-1} \Delta H^{(i)}_{\mathrm{sell}}}{V_t}
    < S^{(j)}_{\mathrm{distress}}.
  \end{equation}
  The cascade is self-limiting iff there exists $j^* < F$ such that the
  post-liquidation price exceeds all remaining thresholds.
\end{theorem}

\begin{proof}
  By induction on the ordered distress-threshold sequence.
\end{proof}

Strategy's distress threshold: $S_{\mathrm{distress}} = (D_t + P_t^{\mathrm{liq}})/H_t
= \$18.45$B$/761{,}068 \approx \$24{,}200$ (full stack).  This is well below
$S_0 = \$73{,}909$, implying substantial buffer.

%-----------------------------------------------------------------------------
\subsubsection{Asymptotic Stability}
%-----------------------------------------------------------------------------

Modelling volume growth as $V_t = V_0 e^{g_V t}$, the time-varying
amplification factor is $R_t = 1/(1 - G_0 e^{(g_E - 2g_S - g_V)t})$.

\begin{corollary}[Asymptotic stability]\label{cor:asymptotic}
  If $2g_S + g_V > g_E$, then $G_t \to 0$ and $R_t \to 1$: the reflexive
  loop vanishes as the BTC market grows relative to any single buyer.
\end{corollary}
